% Pursuing Stacks (cotes 134-2 to 134-8): archive cote 134-8, site page 21 (PDF page 22), typescript page 575.
% typescript. Draft compared: ps7.tex:968-1030.
% AI-transcribed by Claude Opus 5.5 (claude-opus-5-5, adaptive thinking, effort high), Message Batches API, batch ps4 (msgbatch_012Aoc1NuskATPsyBeKSYK7J), 2026-11-12; reviewed against both supplied scan halves; see CHECK_DECISIONS.tsv.
\scanpage{134-8/21}
\mspage{575}

This is in sharp contrast with the more naïve notions of ``cohomological''
properness and smoothness via base change properties, which are inter-%
changed by duality (which was part of yesterday's perplexities,\over{n}{\unsure{N|n}}ow
straightening out ...).

I feel I should be a little more outspoken about the relevant
notion of ``proper maps'' between ordered sets, which should be the alge-%
braic expression of the geometric notion alluded to above\hand{,} of a map
between conically stratified spaces being ``compatible with the conical
stratifications''. In order for the corresponding direct image functor
for sheaves to take sheaves compatible with the stratification above to
sheaves of s ame type below, we'll have to insist that the image by f
of a strata above should be a whole strata below. Now it is clear that
any map
\[ (9) \qquad f: I \hand{\longrightarrow} J \]
between ordered sets \xout{defines} takes flags into flags, and hence induces
a map
\[ (10) \qquad \hand{|}f\hand{|}:\hand{|}I\hand{|}\hand{\to}\hand{|}J\hand{|} \]
between the geometric realizations, compatible with the barycentric
triangulations and \deleted{hence} \inserted{thus} taking simplices into simplices. But even
when I and J are finite (hence ``proper''), the latter map does not al-%
ways satisfy the condition above\over{,}{.} \xout{even when} Thus when I is reduced to
just one point, hence (9) is defined \xout{\illegible{1}i} by the image $\over{\illegible{1}}{y}\hand{\in} J$ of the lat-%
ter, the corresponding map (10) maps the unique point of $\hand{|}I\hand{|}$ to the
barycenter of the strat\over{a}{um} $C_{\over{\illegible{1}}{y}}$ in $\hand{|}J\hand{|}$ , which is a strat\over{a}{um} of $\hand{|}J\hand{|}$ iff
\xout{\illegible{1} is minim} $C_{\over{\illegible{1}}{y}}$ is a minimal strata, i.e. $\over{\illegible{1}}{y}$ is a \uline{maximal} element in J.
A natural algebraic condition to impose upon f, in order to insure
that $\hand{|}f\hand{|}$ satisfy the simple geometric condition above, is that for
any x in I, and any y' in J such that
\[ y' \hand{>} y \overset{\text{def}}{=} f(x) \qquad , \]
there should exist an x' in I satisfying
\[ \hand{(11)} \qquad x' \hand{>} x \quad \text{and} \quad f(x') = y' \quad . \]
This condition strongly reminds of the valuative criterion for proper-%
ness in the context of schemes, where the relation $y' \hand{\geq} y$ \inserted[typed]{or $y\hand{\to} y'$}\hand{,} say\hand{,} should
be interpreted as meaning that y' is a \uline{specialization} of \hand{$y$} . However,
in the valuative criterion for properness (for a map of preschemes of
finite type over a noetherian prescheme say), if one wants actual pro\unsure{p}-%

\margin{21}

% Visual collation: both supplied image halves inspected; see CHECK_DECISIONS_134_7_8.tsv.
